\section{Three data sets}
\label{sec:realdata}

\subsection{The fishery data}
\label{sec:fishery}

The fishery data \citep{riani2008fitting}, distributed with FSDA, record $677$ monthly import
flows of one fishery product into the European Union from a third country: the quantity and
the value of each flow. They are the standard example for TCLUST-REG and the forward search
\citep{perrotta2009new,cerioli2014robust}: flows priced differently form lines through the
origin, and a number of flows fit none of them. We work on the logarithmic scale, where value
$=$ price $\times$ quantity becomes a line of slope one whose intercept is the log price, the
spread of the errors no longer grows with the quantity, and the concentration of small flows
near the origin, which \citet{cerioli2014robust} handle by thinning, disappears. Twenty-six
flows with zero quantity or value have no logarithm and are left out, which leaves $651$. The
published analyses \citep{perrotta2009new,cerioli2014robust,torti2019assessing} keep them and
work on the raw scale, where the heteroscedasticity and the dense region near the origin are
part of the problem; the analysis below therefore answers a different question and is not a
comparison with their results. We fit $K=2$ groups with the constants of Section~\ref{sec:constants} and
$m=8$.

Supplement Table~S42 and Figure~S1 show the results. ESF alone finds two lines
of slope close to one, with intercepts $2.57$ and $2.85$, that is prices of about $13$ and
$17$, and flags $11\%$ of the flows, about $80\%$ of them priced below both lines. The recovery
step finds that these cheap flows lie close to a line of their own: for $19$ of $20$ seeds it
replaces the line at $2.57$ by a line with intercept $2.16$, a price of about $9$, keeps one
line through the two dearer price levels (intercept $2.78$), and flags $3\%$ of the flows.

TCLUST-REG gives the same two answers, but which one depends on the trimming level: at levels
up to $0.08$ it treats the cheap flows as a group, with intercepts between $2.2$ and $2.3$, as
ESF does; from level $0.20$ on it separates the two dearer price levels, as ESF alone does,
and its estimates then vary from one random start to another (Supplement, Section~S5). With the long search, reweighting and the recovery step,
level $0.25$ returns the cheap group and the dearer line (intercepts $2.15$ to $2.16$ and
$2.78$, $3\%$ of the flows flagged) from all ten starts, and level $0.40$ from three of ten;
without the recovery step both levels separate the two dearer price levels.

The data suggest three price levels. With $K=3$ and $m=16$, ESF returns lines with intercepts
$1.90$ to $1.91$, $2.50$ to $2.54$ and $2.82$ to $2.83$, slopes between $0.94$ and $1.00$, and
$3\%$ of the flows flagged, for $12$ of $20$ seeds; the other seeds give other fits. The
best-of-five rule of Section~\ref{sec:constants} returned the cheap group and the dearer line
with $K=2$, and the three price levels with $K=3$, every time over twenty disjoint sets of
seeds (Supplement, Section~S5). The Bayesian information criterion of the same model
prefers $K=2$ (Section~\ref{sec:discussion}); whether the cheap flows are a price level or
anomalous declarations is a question for the subject-matter expert.


\subsection{The tone perception data}
\label{sec:tone}

The tone perception data \citep{deveaux1989mixtures}, distributed with the R package
\texttt{mixtools} \citep{benaglia2009mixtools}, record $150$ trials in which a musician tuned a
tone to the octave above a fundamental with stretched overtones; the tuned ratio lies either
near $2$ or near the stretch ratio, so the data form a nearly flat line and a line of slope
one. \citet{bai2012robust} added ten identical points at $(0,4)$, far
from both lines and outside the range of the covariate. We fit $K=2$ with the constants of
Section~\ref{sec:constants} and $m=8$, to the original data and with the ten points added
(Supplement, Table~S35). On the original data every method finds the two lines, except the
trimmed cluster-weighted model at level $0.25$, which returns other fits from two of its ten
starts. With the ten points added, $6\%$ of the enlarged sample, one of the two lines is lost
from every start by TCLUST-REG at levels $0.02$ and $0.05$, by the trimmed cluster-weighted
model at level $0.05$ and by the Laplace mixture; the trimmed cluster-weighted model at levels
$0.10$ and $0.25$ loses it from three and six of its ten starts, and the noise-component
mixture for $7$ of $20$ seeds. ESF returns the same two
lines, to within $0.01$, for every seed, as do CTLE, the contaminated normal mixture and
TCLUST-REG from level $0.10$ on, with either search; the recovery step found no group among the
flagged trials. The flagged fraction of ESF, $0.25$ to $0.27$ on the original data, is not a
count of gross outliers: a tight core of trials sets the robust scale of the rising line, so
$18$ to $20$ trials of that line are flagged by their residuals, and the covariate screen flags
$14$ more (Supplement, Section~S5); TCLUST-REG with reweighting flags $19\%$. The fitted
lines are not affected.

\subsection{Taxi fares from an airport}
\label{sec:taxi}

The third example has a known answer. The New York City Taxi and Limousine Commission publishes
every yellow-taxi trip \citep{tlc2024trip}; we took the $157{,}264$ trips of March 2024 that
started at John F.\ Kennedy airport and carry a rate code. Their fare, before tolls, surcharges
and tips, follows one of two tariffs: a flat fare of $70$ dollars to Manhattan (rate code 2
with a valid record, $53\%$ of the trips), or the meter, which charges by distance and, in slow
traffic, by time (rate code 1, $38\%$); the remaining $9.6\%$ use other tariffs or are invalid
records (Supplement, Section~S5) and are treated as contamination. We regressed the fare on the
distance and the duration, $d=3$, with $K=2$ and $m=10$; nothing was removed, and the rate
code gives the true group of every clean trip. Twenty random samples of $2{,}000$ trips were
fitted by every method of Section~\ref{sec:simulation} except sequential RANSAC, which needs
the noise scale. The flat-fare group lies exactly on its line, so more than half of the units
have residual zero and the scale of the scoring step, a median of squared residuals, is zero;
the guard mentioned in Section~\ref{sec:design} takes the median over the units that a
replicate does not fit exactly.

Supplement Table~S43 gives the results. ESF found both tariffs in all twenty samples: the flat
fare exactly, and a metered line with median intercept $2.63$ dollars, $3.04$ dollars per mile
and $0.33$ dollars per minute, against $2.40$, $3.06$ and $0.33$ for the fit to all metered
trips. Its accuracy of $0.978$ is the ceiling that the best methods share: TCLUST-REG and the
trimmed cluster-weighted model at levels $0.10$ to $0.40$, the bisquare, contaminated normal and
noise-component mixtures all reach $0.975$ to $0.979$, as does TCLUST-REG with the long search,
reweighting and the recovery step at both levels, with smaller coefficient errors ($0.10$ and
$0.14$ against $0.23$); the recovery step changed none of its fits. TLE at levels $0.10$ and
$0.15$ and CTLE stopped with an error on most samples, presumably because the flat-fare group
has no variance; the example shows that the methods hold on a group without noise, not that
one is better than another. ESF took a median of $31$ seconds
per sample, most of it in the minimum covariance determinant of the covariate screen, for which
we used a general-purpose implementation (Supplement, Table~S43). Its flagged fraction, $0.23$, is more than twice the contamination: the metered fares depend on the
time spent in slow traffic, not on the duration, and their residuals have heavy tails, so
$16\%$ of the clean trips are flagged, more than half of them by the covariate screen, which
reads the long, skewed trips as outlying covariates; $89\%$ of the contaminated trips are
flagged. TCLUST-REG with reweighting, which has no covariate screen, flags $0.14$. The
best-of-five rule of Section~\ref{sec:constants} failed here once, because with a group that
has no noise the likelihood rewards any line through tied points.
